% =============================================================================
% CompetitiveExam/CS301/07-linked-lists-doubly-applications-answers.tex
% Answer key for 07-linked-lists-doubly-applications-questions.tex.
%
% Real GATE PYQs (Q1--Q5) are sourced from the local GATEOverflow Volume 2
% PYQ book (CompetitiveExam/Books/index.md, Tier-1 availability: Linked List
% cluster PDF pp.232-238, Queue cluster pp.240-243), which authenticates each
% question (year + question number printed on the page). The book's answer
% keys, assumed collapsed by the registry, are in fact recoverable via
% pdftotext from the volume's "Answer Keys" section; those keys give:
%   3.10.6  (1997 Q1.4)  = C
%   3.10.12 (2004 Q36)   = A
%   3.12.10 (2017 S2 Q13)= B
%   3.12.11 (2018 Q3)    = B
%   3.10.20 (2023 Q3)    = D
% Independent web corroboration obtained this session (2026-10-03):
%   - Q1  stem + options matched on PracticePaper's GATE CSE 1997 paper page.
%   - Q4  all four options + answer B on GeeksforGeeks' GATE-CS-2018
%         linked-list question page.
%   - Q5  all four options on GeeksforGeeks' GATE-CS-2023 question page;
%         answer D corroborated by the GATEQA listing ("Answer: D").
%   (Search engines were partly bot-blocked and no OCR was available; Q2/Q3
%    rest on the local Tier-1 book text + its answer key.)
% Original questions (Q6--Q10) have answers derived directly from the Week 7
% deck algorithms and checked by hand this session.
% =============================================================================

\documentclass[11pt]{article}
\usepackage[margin=1in]{geometry}
% =============================================================================
% Preambles/header.tex — shared LaTeX/Beamer preamble for this project
%
% Use from a Beamer lecture by prepending:
%     \input{header}
% and compiling with TEXINPUTS=../Preambles:$TEXINPUTS (the /compile-latex
% skill does this automatically).
%
% PALETTE CONTRACT. Color names defined here MUST match the names in
% Quarto/theme-template.scss so Beamer and Quarto renderings use the same
% palette. The scripts/check-palette-sync.sh script enforces this.
%
% To customize: edit the HEX values below AND the matching $variable in
% Quarto/theme-template.scss, then run scripts/check-palette-sync.sh.
% =============================================================================

% -----------------------------------------------------------------------------
% Engine detection + encoding
%
% Our /compile-latex skill uses XeLaTeX, where inputenc is unnecessary and
% can produce encoding warnings. Only load inputenc under pdfTeX.
% -----------------------------------------------------------------------------
\usepackage{iftex}
\ifPDFTeX
  \usepackage[utf8]{inputenc}
\fi

\usepackage{amsmath, amssymb, amsthm}
\usepackage{booktabs}
\usepackage{graphicx}

% -----------------------------------------------------------------------------
% PALETTE — mirrors Quarto/theme-template.scss variable names.
% Load xcolor and define colors BEFORE anything that references them
% (notably hyperref, hypersetup, and the Beamer color assignments).
% -----------------------------------------------------------------------------
\usepackage{xcolor}

\definecolor{primary-blue}{HTML}{012169}      % headings, accents
\definecolor{primary-gold}{HTML}{B9975B}      % emphasis, borders
\definecolor{highlight-yellow}{HTML}{F2A900}  % markers, alerts
\definecolor{light-bg}{HTML}{E8EDF5}          % title / section backgrounds
\definecolor{jet}{HTML}{1A1A1A}               % body text

% Semantic colors (used in diagrams and annotations)
\definecolor{positive}{HTML}{15803D}          % good / observed / identified
\definecolor{negative}{HTML}{B91C1C}          % bad / problematic / bias
\definecolor{neutral}{HTML}{525252}           % reference / context

% Highlight accents (match .hi-* classes in the SCSS)
\definecolor{hi-slate}{HTML}{314F4F}
\definecolor{hi-green}{HTML}{2E8B57}
\definecolor{hi-red}{HTML}{C41E3A}

% Semantic aliases so skill templates and snippets can write
% \textcolor{accent}{...} without hard-coding "primary-blue".
\colorlet{accent}{primary-blue}
\colorlet{accent2}{primary-gold}

% -----------------------------------------------------------------------------
% Hyperref — loaded AFTER xcolor + palette so link colors resolve.
% -----------------------------------------------------------------------------
\usepackage{hyperref}
\hypersetup{colorlinks=true, linkcolor=primary-blue, urlcolor=primary-blue,
            citecolor=primary-blue, pdfborder={0 0 0}}

% -----------------------------------------------------------------------------
% Course code — set once per deck (after \input{header}, before \begin{document})
% via \coursecode{CS401}. Shown in the footer on every slide so a deck's course
% is identifiable without opening the title page. Empty by default.
% -----------------------------------------------------------------------------
\makeatletter
\newcommand{\coursecode}[1]{\gdef\@coursecodetext{#1}}
\gdef\@coursecodetext{}
\makeatother

% -----------------------------------------------------------------------------
% Beamer theme assignments (only apply under Beamer).
%
% We detect Beamer via \@ifclassloaded{beamer}, which is the documented,
% non-internal way. This must be wrapped in \makeatletter/\makeatother
% because \@ifclassloaded contains an @-catcode character.
% -----------------------------------------------------------------------------
\makeatletter
\@ifclassloaded{beamer}{%
  \usefonttheme{professionalfonts}
  \setbeamercolor{structure}{fg=primary-blue}
  \setbeamercolor{frametitle}{fg=primary-blue}
  \setbeamercolor{title}{fg=primary-blue}
  \setbeamercolor{subtitle}{fg=primary-gold}
  \setbeamercolor{author}{fg=primary-blue}
  \setbeamercolor{institute}{fg=neutral}
  \setbeamercolor{date}{fg=primary-gold}
  \setbeamercolor{itemize item}{fg=primary-blue}
  \setbeamercolor{itemize subitem}{fg=primary-gold}
  \setbeamercolor{alerted text}{fg=primary-gold}
  \setbeamercolor{block title}{fg=white, bg=primary-blue}
  \setbeamercolor{block body}{bg=light-bg}
  % Semantic block variants: block=definition/concept (blue), exampleblock=
  % worked example (green/positive), alertblock=key takeaway or caution (gold).
  % Keep to <=2 colored boxes per slide across ALL three variants (INV-7).
  \setbeamercolor{block title example}{fg=white, bg=positive}
  \setbeamercolor{block body example}{bg=light-bg}
  \setbeamercolor{block title alerted}{fg=white, bg=primary-gold}
  \setbeamercolor{block body alerted}{bg=light-bg}

  % Minimal footer: course code (left) + page X / N (right), no navigation clutter
  \setbeamertemplate{navigation symbols}{}
  \setbeamertemplate{footline}{%
    \color{neutral}\footnotesize\hspace{0.7em}\@coursecodetext\hfill
    \insertframenumber/\inserttotalframenumber\hspace{0.7em}\vspace{0.3em}}
}{}
\makeatother

% -----------------------------------------------------------------------------
% TikZ setup — libraries the snippet gallery uses
% -----------------------------------------------------------------------------
\usepackage{tikz}
\usetikzlibrary{arrows.meta, positioning, calc, decorations.pathreplacing,
                fit, shapes.geometric, backgrounds}

% Shared TikZ styles — snippets and hand-written diagrams can pick these up
% via `\begin{tikzpicture}[every node/.style={...}]` or by naming specific
% styles (e.g., `\node[dag-node] (X) at (0,0) {X};`).
\tikzset{
  % Node styles for causal DAGs and similar
  dag-node/.style = {
    draw=primary-blue, fill=light-bg, rounded corners=2pt,
    minimum width=1.2cm, minimum height=0.9cm, align=center, font=\small
  },
  decision-node/.style = {
    draw=primary-gold, fill=white, diamond, aspect=1.8,
    minimum width=1.8cm, minimum height=1.2cm, align=center, font=\small,
    inner sep=1pt
  },
  % Edge styles — observed vs counterfactual
  observed-edge/.style = {-{Latex[length=2.5mm]}, thick, draw=jet},
  counterfactual-edge/.style = {-{Latex[length=2.5mm]}, thick, draw=neutral, dashed},
  confound-edge/.style = {-{Latex[length=2.5mm]}, thick, draw=negative, dashed},
  % Data-point styles for plots
  observed-dot/.style = {fill=primary-blue, draw=primary-blue, circle, inner sep=1.2pt},
  counterfactual-dot/.style = {fill=white, draw=neutral, circle, inner sep=1.2pt}
}

% -----------------------------------------------------------------------------
% Convenience macros
% -----------------------------------------------------------------------------
\newcommand{\muted}[1]{\textcolor{neutral}{#1}}
\newcommand{\key}[1]{\textcolor{primary-gold}{\textbf{#1}}}
\newcommand{\good}[1]{\textcolor{positive}{#1}}
\newcommand{\bad}[1]{\textcolor{negative}{#1}}

% Standout frame macro: full-bleed dark background for section transitions.
% Beamer-only (uses \begin{frame}/\setbeamercolor/\usebeamerfont) -- do not
% call from an article-class file (e.g. Notes/<CODE>/*-notes.tex). \muted,
% \key, \good, \bad, and \coursecode above are plain xcolor/gdef and are
% safe to use from either class; verified when Notes/article-class support
% was added.
% Expand: \transitionslide{Your transition title}
\newcommand{\transitionslide}[1]{%
  \begingroup
  \setbeamercolor{background canvas}{bg=primary-blue}
  \begin{frame}[plain]
    \centering
    \vfill
    {\usebeamerfont{title}\color{white}\Large #1\par}
    \vfill
  \end{frame}
  \endgroup
}

% Section divider frame (Beamer-only), an alternative to \transitionslide with a
% darker two-line style: near-black (jet) background, a small white label line
% above a larger gold title, and a thin gold rule along the bottom edge.
% Expand: \sectiondivider{Section 2}{Pointers}
\newcommand{\sectiondivider}[2]{%
  \begingroup
  \setbeamercolor{background canvas}{bg=jet}
  \begin{frame}[plain]
    \vspace*{3.0cm}
    {\color{white}\ttfamily\large #1\par}
    \vspace{0.5cm}
    {\color{primary-gold}\LARGE #2\par}
    \begin{tikzpicture}[remember picture, overlay]
      \draw[primary-gold, line width=2.5pt]
        ($(current page.south west)+(0,0.1)$) -- ($(current page.south east)+(0,0.1)$);
    \end{tikzpicture}
  \end{frame}
  \endgroup
}

% -----------------------------------------------------------------------------
% Channel watermark — "Concepts That Click" (YouTube).
%
% Article-class files (CompetitiveExam Q/A sets, Notes, Assignments) get a
% light diagonal draft watermark on every page plus a centered footer naming
% the channel. Beamer decks get a subtle TikZ background watermark instead
% (the footline already carries the course code). Centralized here so every
% MCQ PDF is branded without per-file edits.
% -----------------------------------------------------------------------------
\makeatletter
\@ifclassloaded{article}{%
  \usepackage{draftwatermark}
  \SetWatermarkText{Concepts That Click}
  \SetWatermarkScale{1}
  \SetWatermarkFontSize{2cm}
  \SetWatermarkLightness{0.86}
  \SetWatermarkAngle{45}
  \usepackage{fancyhdr}
  \pagestyle{fancy}
  \fancyhf{}
  \fancyfoot[C]{\footnotesize\textcolor{neutral}{Concepts That Click~|~YouTube: \href{https://www.youtube.com/@concepts.thatclick}{@concepts.thatclick}}}
  \renewcommand{\headrulewidth}{0pt}
  \renewcommand{\footrulewidth}{0pt}
  \fancypagestyle{plain}{%
    \fancyhf{}%
    \fancyfoot[C]{\footnotesize\textcolor{neutral}{Concepts That Click~|~YouTube: \href{https://www.youtube.com/@concepts.thatclick}{@concepts.thatclick}}}%
    \renewcommand{\headrulewidth}{0pt}%
    \renewcommand{\footrulewidth}{0pt}%
  }
}{}
\@ifclassloaded{beamer}{%
  \setbeamertemplate{background}{%
    \begin{tikzpicture}[remember picture, overlay]
      \node[rotate=45, opacity=0.055, align=center]
        at (current page.center)
        {\Huge\textcolor{neutral}{Concepts That Click}};
    \end{tikzpicture}%
  }
}{}
\makeatother 
\coursecode{CS301}
\renewcommand{\thesection}{7.\arabic{section}}

\title{Competitive Exam Practice 7 --- Answer Key\\\large Doubly \& Circular Lists, List-Backed Stack and Queue, and Polynomials}
\author{Auqib Hamid Lone\\[0.3em]
  \emph{Department of Computer Science \& Engineering}, \emph{GCET Ganderbal}}
\date{\today}

\begin{document}
\maketitle

\begin{enumerate}
\item[\textbf{Q1}] \textbf{(C) Circular doubly linked list.} To concatenate two lists in $O(1)$ you must reach both ends without a walk and be able to rewire the four boundary links; a circular doubly linked list stores a pointer to the last node whose \texttt{next} is the first node and whose \texttt{prev} is the previous last node, so both joins are constant-time splices. A singly or ordinary doubly linked list needs an $O(n)$ walk to reach the tail, and an array needs $O(n)$ copying/reallocation. (Local book answer key 3.10.6 = C; stem + options also on PracticePaper's GATE CSE 1997 page.)

\item[\textbf{Q2}] \textbf{(A) rear node.} Point \texttt{p} at the rear node. Then the front node is \texttt{p->next} (one link, $O(1)$), so \texttt{dequeue} removes at the front and \texttt{enqueue} inserts just after \texttt{p} and advances \texttt{p} to the new node --- both constant time. Pointing at the front instead makes \texttt{enqueue} at the rear need an $O(n)$ walk to find the tail; two pointers would work but the question allows a single one, so (A) is the answer and (C) is false. (Local book answer key 3.10.12 = A.)

\item[\textbf{Q3}] \textbf{(B) (II) only.} For a singly linked list to behave as a circular queue with $O(1)$ insert/delete, the ring must be closed from the rear back to the front: the rear node's \texttt{next} points to the front node, so appending after \texttt{REAR} and advancing \texttt{REAR} keep \texttt{FRONT} reachable. Statement (I) points the front node forward at the rear, which does not make the rear reach the front and is not required. (Local book answer key 3.12.10 = B.)

\item[\textbf{Q4}] \textbf{(B) $\Theta(1)$, $\Theta(n)$.} Inserting at the head is a constant-time front insert ($O(1)$). Deleting at the tail of a \emph{non-circular} singly linked list requires the tail's predecessor to redirect the tail's \texttt{next} to \texttt{NULL}; with only a tail pointer there is no backward link, so the predecessor must be found by walking from the head --- $\Theta(n)$ worst case. Hence \texttt{enqueue} is $\Theta(1)$ and \texttt{dequeue} is $\Theta(n)$. (Local book answer key 3.12.11 = B; all four options on GeeksforGeeks.)

\item[\textbf{Q5}] \textbf{(D) \texttt{SLLdel} is $O(n)$ and \texttt{DLLdel} is $O(1)$.} In a doubly linked list a node reaches both neighbours through \texttt{prev} and \texttt{next}, so \texttt{DLLdel} unlinks it in $O(1)$ given a pointer to it. In a singly linked list the same node cannot reach its predecessor, so unless it is the first node (or the copy-next trick is used, which fails for the last node) deleting it requires an $O(n)$ walk to find the node whose \texttt{next} points at it. (Local book answer key 3.10.20 = D; all four options on GeeksforGeeks; GATEQA lists answer D.)

\item[\textbf{Q6}] \textbf{(A) The list becomes \texttt{head} $\leftrightarrow$ 7 $\leftrightarrow$ 42; the two boundary links are rewired in $O(1)$ and the delete is $O(1)$ once \texttt{p} is known.} With \texttt{p} at the 19-node, \texttt{p->prev} is the 7-node and \texttt{p->next} is the 42-node, both in hand: set \texttt{7->next = p->next} (42) and \texttt{42->prev = p->prev} (7), then \texttt{free(p)}. (B) wrongly charges an $O(n)$ search the backward link eliminates; (C) denies the bypass ever happens; (D) updates only one link, leaving 42's \texttt{prev} dangling --- the classic half-invariant bug. Confirmed by trace (2026-10-03).

\item[\textbf{Q7}] \textbf{(A).} Wire the new node forward and backward while \texttt{p->next} still names the old successor (14), repair 14's backward link, and only then overwrite \texttt{p->next} --- i.e.\ \texttt{newNode->next = p->next}, then \texttt{newNode->prev = p}, then \texttt{p->next->prev = newNode}, and finally \texttt{p->next = newNode}. (B) overwrites \texttt{p->next} first, so \texttt{newNode->next} becomes a self-loop and 14 is lost. (C) and (D) overwrite \texttt{p->next} too early, so the later \texttt{p->next->prev = newNode} sets \texttt{newNode->prev} to itself and leaves 14's \texttt{prev} pointing at 8 --- a broken two-link invariant. Confirmed by trace (2026-10-03).

\item[\textbf{Q8}] \textbf{(B) It loops forever, cycling p, q, r, \dots} In a ring no node's \texttt{next} is \texttt{NULL}, so the test \texttt{ptr != NULL} is never false and the walk never terminates. The correct traversal tests \texttt{ptr != head} (preferably with a \texttt{do--while} so the head is visited once). (A) would require a \texttt{NULL} terminator that a circle does not have; (C) misreads the loop condition as checked once; (D) is not what the loop does. Confirmed by trace (2026-10-03).

\item[\textbf{Q9}] \textbf{(A) $A + B = 10x^{3} + 5x^{2} + 9$, one merge pass, $O(n+m)$; the cancelled $x$ term is not stored.} Merge by descending exponent: exp~3 gives $6 + 4 = 10$; exp~2 exists only in $B$, so $5x^{2}$ copies through; exp~1 gives $2 + (-2) = 0$, and a zero-coefficient term is never emitted; exp~0 gives 9. Each source node is visited once, so the cost is linear. (B) stores a forbidden zero term; (C) confuses addition with the $O(n \cdot m)$ multiplication; (D) wrongly keeps $4x$. Confirmed by trace (2026-10-03).

\item[\textbf{Q10}] \textbf{(A) \texttt{push} front-inserts 28 giving \texttt{head} $\to$ 28 $\to$ 7 $\to$ 14 $\to$ 21; \texttt{pop()} returns 28, leaving top 7; both $O(1)$ worst case.} The stack top is \texttt{head}, so \texttt{push} is Week 6's front-insert and \texttt{pop} is front-delete --- each rewrites \texttt{head} and one link, both constant time. (B) appends at the tail and mis-states which end pops; (C) mislabels the LIFO discipline (a stack returns the newest, 28, not the oldest); (D) is false: no walk is needed at \texttt{head}. Confirmed by trace (2026-10-03).
\end{enumerate}

\end{document}
